% GENERATED FILE. Edit canonical lessons, then run npm run generate:books.
% canonical-source-hash: fnv1a64:3baea3829060e4d5
% canonical-lessons: FA-W22-re, FA-W22-ze, FA-W22-dal

\chapter{Letters That Never Join Forward}
\label{ch:fa-letters-never-join-forward}
\hlchaptermodality{22}

\begin{chapteropening}
\textbf{By the end of this chapter:} \emph{I can write re, ze and dal, and find each one in a word I already say.}

Complete the last lesson of chapter 22: i can write re, ze and dal, and find each one in a word I already say.
\end{chapteropening}

\section[re]{\fa{ر} (re) — the letter that ends \fa{چطور}}

\label{lesson:FA-W22-re}



\begin{quote}
Before the new one: write \fa{ه} and \fa{و} once each, and name them.

Now say \textbf{\fa{چطور}} (\emph{chetor}), \textquotedblleft{}how\textquotedblright{}. You have read it for a long time. This lesson takes one letter out of it and puts it in your hand.
\end{quote}

\subsection*{Script you'll notice: \fa{ر}}
\textbf{\fa{ر}} — \emph{re}, the sound \emph{r}.

It ends \textbf{\fa{چطور}} (\emph{chetor}), \textquotedblleft{}how\textquotedblright{}, from \textquotedblleft{}how are you?\textquotedblright{}. \textbf{\fa{ر}} is a \textbf{non-connector}, like \textbf{\fa{ا}} and \textbf{\fa{و}}, which you already write: it joins the letter before it but never the one after, so a small gap always follows it.

\subsection*{Writing: \fa{ر}}
\hlblockfigure{\detokenize{figures/FA-W22-re-filmstrip.pdf}}{How \fa{ر} is written, stroke by stroke}

\begin{itemize}
\raggedright
  \item \textbf{1.} start here: begin at the upper tip and descend through the short stroke
  \item \textbf{2.} without lifting, sweep left through the lower curve — and only now lift
\end{itemize}

\textbf{Pen lifts: 0.} The pen never leaves the paper.

\begin{quote}
Stroke order is one attested teaching order; hands and teachers vary. Source: Persian Online, How to Write Persian Characters, \fa{ر} demonstration at 01:10–01:12 (Liberal Arts Instructional Technology Services, Univ. of Texas at Austin).
\end{quote}

\subsection*{Your turn}
\begin{itemize}
\raggedright
  \item \emph{Look:} at \textbf{\fa{چطور}} and point at \fa{ر} inside it
  \item \emph{Trace it:} \fa{ر} three times, saying \emph{r} as you finish each one
  \item \emph{Write it:} \fa{و}, then \fa{ر}, from memory
  \item \emph{Read it:} \textbf{\fa{چطور}} aloud once more
\end{itemize}

\begin{tcolorbox}[breakable,colback=teal!4,colframe=teal!35!black,title={Before you move on}]
Which letter is this — \fa{ر}? (\textbf{re}, \emph{r}.) Which word did it come from? (\textbf{\fa{چطور}}, \emph{chetor}, \textquotedblleft{}how\textquotedblright{}.)
\end{tcolorbox}
\section[ze]{\fa{ز} (ze) — the letter that opens \fa{زبان}}

\label{lesson:FA-W22-ze}



\begin{quote}
Before the new one: write \fa{و} and \fa{ر} once each, and name them.

Now say \textbf{\fa{زبان}} (\emph{zabân}), \textquotedblleft{}language\textquotedblright{}. You have read it for a long time. This lesson takes one letter out of it and puts it in your hand.
\end{quote}

\subsection*{Script you'll notice: \fa{ز}}
\textbf{\fa{ز}} — \emph{ze}, the sound \emph{z}.

It opens \textbf{\fa{زبان}} (\emph{zabân}), \textquotedblleft{}language\textquotedblright{}, the thing you are learning. It is \textbf{\fa{ر}}, from the last lesson, with \textbf{one dot above}, and like \textbf{\fa{ر}} it never joins the letter after it.

\subsection*{Writing: \fa{ز}}
\hlblockfigure{\detokenize{figures/FA-W22-ze-filmstrip.pdf}}{How \fa{ز} is written, stroke by stroke}

\begin{itemize}
\raggedright
  \item \textbf{1.} start here: begin at the upper tip and descend through the short stroke
  \item \textbf{2.} without lifting, sweep left through the lower curve
  \item \textbf{3.} lift once, then place the dot above — and only now lift
\end{itemize}

\textbf{Pen lifts: 1.}

\begin{quote}
Stroke order is one attested teaching order; hands and teachers vary. Source: Persian Online, How to Write Persian Characters, \fa{ز} demonstration at 01:13–01:16 (Liberal Arts Instructional Technology Services, Univ. of Texas at Austin).
\end{quote}

\subsection*{Your turn}
\begin{itemize}
\raggedright
  \item \emph{Look:} at \textbf{\fa{زبان}} and point at \fa{ز} inside it
  \item \emph{Trace it:} \fa{ز} three times, saying \emph{z} as you finish each one
  \item \emph{Write it:} \fa{ر}, then \fa{ز}, from memory
  \item \emph{Read it:} \textbf{\fa{زبان}} aloud once more
\end{itemize}

\begin{tcolorbox}[breakable,colback=teal!4,colframe=teal!35!black,title={Before you move on}]
Which letter is this — \fa{ز}? (\textbf{ze}, \emph{z}.) Which word did it come from? (\textbf{\fa{زبان}}, \emph{zabân}, \textquotedblleft{}language\textquotedblright{}.)
\end{tcolorbox}
\section[dâl]{\fa{د} (dâl) — the letter inside \fa{خدا}}

\label{lesson:FA-W22-dal}



\begin{quote}
Before the new one: write \fa{ر} and \fa{ز} once each, and name them.

Now say \textbf{\fa{خدا}} (\emph{khodâ}), \textquotedblleft{}God\textquotedblright{}. You have read it for a long time. This lesson takes one letter out of it and puts it in your hand.
\end{quote}

\subsection*{Script you'll notice: \fa{د}}
\textbf{\fa{د}} — \emph{dâl}, the sound \emph{d}.

It is inside \textbf{\fa{خدا}} (\emph{khodâ}), the first half of the farewell \textbf{\fa{خدا} \fa{حافظ}}. It is the third letter that never joins forward: after \textbf{\fa{د}}, as after \textbf{\fa{ر}} and \textbf{\fa{ا}}, the pen lifts and the next letter starts fresh.

\subsection*{Writing: \fa{د}}
\hlblockfigure{\detokenize{figures/FA-W22-dal-filmstrip.pdf}}{How \fa{د} is written, stroke by stroke}

\begin{itemize}
\raggedright
  \item \textbf{1.} start here: begin at the upper tip and descend through the folded shoulder
  \item \textbf{2.} turn left along the baseline without lifting — and only now lift
\end{itemize}

\textbf{Pen lifts: 0.} The pen never leaves the paper.

\begin{quote}
Stroke order is one attested teaching order; hands and teachers vary. Source: Persian Online, How to Write Persian Characters, \fa{د} demonstration at 01:04–01:06 (Liberal Arts Instructional Technology Services, Univ. of Texas at Austin).
\end{quote}

\subsection*{Your turn}
\begin{itemize}
\raggedright
  \item \emph{Look:} at \textbf{\fa{خدا}} and point at \fa{د} inside it
  \item \emph{Trace it:} \fa{د} three times, saying \emph{d} as you finish each one
  \item \emph{Write it:} \fa{ز}, then \fa{د}, from memory
  \item \emph{Read it:} \textbf{\fa{خدا}} aloud once more
\end{itemize}

\begin{tcolorbox}[breakable,colback=teal!4,colframe=teal!35!black,title={Before you move on}]
Which letter is this — \fa{د}? (\textbf{dâl}, \emph{d}.) Which word did it come from? (\textbf{\fa{خدا}}, \emph{khodâ}, \textquotedblleft{}God\textquotedblright{}.)
\end{tcolorbox}
